%%%%%%%%%%  scan idx 477  |  folio 464  %%%%%%%%%%
\textbf{Remarque 4.9.} Dans le cas \(X\) noethérien, \(Y\) régulier, on peut, en procédant comme en 4.7, décrire la structure d'anneau graduée de \(\operatorname{Gr}^*K^\circ(X')_{\mathbb Q}\) en termes des anneaux gradués \(\operatorname{Gr}^*K^\circ(X)_{\mathbb Q}\), \(\operatorname{Gr}^*K^\circ(Y)_{\mathbb Q}\), des homomorphismes additifs \(i^{\gr}\) et \(i_{\gr}\) déduits de \(i^*\) et \(i_*\), et de l'élément \(c^{\bullet}(N)\) de \(\operatorname{Gr}^*K^\circ(Y)\). De même, on peut décrire à partir des mêmes données les homomorphismes \(f^{\gr},f_{\gr},j^{\gr},j_{\gr}\) définis à partir de \(f^*,f_*,j^*,j_*\) par passage aux gradués associés.

\subsection*{4.10.}

Soit \(X\) un schéma \emph{régulier} ayant un faisceau ample. On sait (IV\qquad) que les groupes de Grothendieck formés avec les \(\mathcal O_X\)-Modules localement libres de type fini, et avec les \(\mathcal O_X\)-Modules cohérents, sont isomorphes. On utilise ce fait pour définir sur \(K^\circ(X)\) une filtration « topologique », en prenant pour \(\Fil^i_{\mathrm{top}}K^\circ(X)\) le sous-groupe engendré par les classes des faisceaux cohérents dont le support est de codimension \(\ge i\) ; on montre facilement que \(\Fil^i_{\mathrm{top}}K^\circ(X)\) est engendré par les classes des faisceaux \(\mathcal O_Z\), où \(Z\) parcourt l'ensemble des sous-schémas fermés intègres de \(X\), de codimension \(\ge i\)(X\qquad).

\textbf{Proposition 4.11.} Soit \(X\) un schéma régulier ayant un faisceau ample. Alors pour tout \(k\ge0\) :
\[
  \Fil_{\mathrm{top}}^kK^\circ(X)_{\mathbb Q}
  \subset
  \Fil_\lambda^kK^\circ(X)_{\mathbb Q}.
\]
\((\mathrm{N.B.}\) L'inclusion inverse, valable sans tensorisation par \(\mathbb Q\), sera prouvée dans X\qquad).

D'après 4.10, il suffit de montrer que si \(Z\) est un sous-schéma fermé intègre de \(X\), de codimension \(k\), alors :
\[
  \operatorname{cl}(\mathcal O_Z)\in\Fil^kK^\circ(X)_{\mathbb Q}.
\]

Pour cela, on procède par récurrence noethérienne sur les sous-schémas fermés intègres de \(X\). Supposons donc la propriété prouvée pour les sous-schémas fermés intègres de \(Z\), et montrons la pour \(Z\). Soir \(J\) l'Idéal de \(\mathcal O_X\)
